
\documentclass{article}
\usepackage{colortbl}
\usepackage{makecell}
\usepackage{multirow}
\usepackage{supertabular}

\begin{document}

\newcounter{utterance}

\centering \large Interaction Transcript for game `dond', experiment `semi\_it', episode 12 with claude{-}sonnet{-}4{-}20250514{-}no{-}reasoning{-}t1.0.
\vspace{24pt}

{ \footnotesize  \setcounter{utterance}{1}
\setlength{\tabcolsep}{0pt}
\begin{supertabular}{c@{$\;$}|p{.15\linewidth}@{}p{.15\linewidth}p{.15\linewidth}p{.15\linewidth}p{.15\linewidth}p{.15\linewidth}}
    \# & $\;$A & \multicolumn{4}{c}{Game Master} & $\;\:$B\\
    \hline

    \theutterance \stepcounter{utterance}  
    & & \multicolumn{4}{p{0.6\linewidth}}{
        \cellcolor[rgb]{0.9,0.9,0.9}{
            \makecell[{{p{\linewidth}}}]{
                \texttt{\tiny{[P1$\langle$GM]}}
                \texttt{State giocando a un gioco di negoziazione in cui dovete accordarvi su come dividere una serie di oggetti tra voi e un altro giocatore.} \\
\\ 
\texttt{Regole:} \\
\texttt{(a) A Lei e all'altro giocatore viene dato un insieme di oggetti. Ognuno di voi riceve anche una funzione di valore segreta, che rappresenta il valore di ciascun tipo di oggetto.} \\
\texttt{(b) Si scambiano messaggi con l'altro giocatore per concordare chi si aggiudica gli oggetti. Potete inviare un massimo di 5 messaggi ciascuno, oppure terminare in anticipo facendo una proposta segreta in qualsiasi momento.} \\
\texttt{(c) A ciascuno di voi viene chiesto di inviare una proposta segreta indicando gli oggetti che desiderate, formattata tra parentesi quadre come segue: "{[}Proposta: <numero> <nome oggetto>, <numero> <nome oggetto>, <...>{]}".} \\
\texttt{(d) Se le vostre proposte sono complementari, cioè ci sono abbastanza oggetti per soddisfare entrambe le proposte, a ciascun giocatore viene assegnato un punteggio basato sulla somma dei valori degli oggetti ricevuti. In caso contrario, entrambi ricevono zero punti.} \\
\texttt{(e) L'obiettivo è massimizzare il proprio punteggio.} \\
\\ 
\texttt{Cominciamo.} \\
\\ 
\texttt{L'insieme degli oggetti disponibili è:} \\
\\ 
\texttt{1 lampada, 3 ditali, 1 corda} \\
\\ 
\texttt{La funzione valore segreta è:} \\
\\ 
\texttt{lampada: 3, ditale: 2, corda: 1} \\
\\ 
\texttt{IMPORTANTE: i vostri messaggi, a meno che non si tratti di una proposta segreta, vengono trasmessi direttamente all'altro giocatore, quindi non includete alcuna risposta alle regole o testo di annuncio del vostro messaggio. Per fare una proposta segreta, utilizzate il formato indicato. Non utilizzare delle parentesi quadre quando si comunica all'altro giocatore, altrimenti verrà interpretata come una proposta segreta.} \\
            }
        }
    }
    & \\ \\

    \theutterance \stepcounter{utterance}  
    & \multicolumn{4}{p{0.6\linewidth}}{
        \cellcolor[rgb]{0.9,0.9,0.9}{
            \makecell[{{p{\linewidth}}}]{
                \texttt{\tiny{[P1$\rangle$GM]}}
                \texttt{Ciao! Vedo che abbiamo a disposizione 1 lampada, 3 ditali e 1 corda da dividere tra noi. Che tipo di oggetti ti interessano di più? Io sono abbastanza flessibile e vorrei trovare una divisione che funzioni bene per entrambi.} \\
            }
        }
    }
    & & \\ \\

    \theutterance \stepcounter{utterance}  
    & & & \multicolumn{2}{p{0.3\linewidth}}{
        \cellcolor[rgb]{0.9,0.9,0.9}{
            \makecell[{{p{\linewidth}}}]{
                \texttt{\tiny{[GM$|$GM]}}
                \texttt{continue} \\
            }
        }
    }
    & & \\ \\

    \theutterance \stepcounter{utterance}  
    & & & \multicolumn{2}{p{0.3\linewidth}}{
        \cellcolor[rgb]{0.9,0.9,0.9}{
            \makecell[{{p{\linewidth}}}]{
                \texttt{\tiny{[GM$|$GM]}}
                \texttt{Ciao! Vedo che abbiamo a disposizione 1 lampada, 3 ditali e 1 corda da dividere tra noi. Che tipo di oggetti ti interessano di più? Io sono abbastanza flessibile e vorrei trovare una divisione che funzioni bene per entrambi.} \\
            }
        }
    }
    & & \\ \\

    \theutterance \stepcounter{utterance}  
    & & \multicolumn{4}{p{0.6\linewidth}}{
        \cellcolor[rgb]{0.9,0.9,0.9}{
            \makecell[{{p{\linewidth}}}]{
                \texttt{\tiny{[GM$\rangle$P2]}}
                \texttt{State giocando a un gioco di negoziazione in cui dovete accordarvi su come dividere una serie di oggetti tra voi e un altro giocatore.} \\
\\ 
\texttt{Regole:} \\
\texttt{(a) A Lei e all'altro giocatore viene dato un insieme di oggetti. Ognuno di voi riceve anche una funzione di valore segreta, che rappresenta il valore di ciascun tipo di oggetto.} \\
\texttt{(b) Si scambiano messaggi con l'altro giocatore per concordare chi si aggiudica gli oggetti. Potete inviare un massimo di 5 messaggi ciascuno, oppure terminare in anticipo facendo una proposta segreta in qualsiasi momento.} \\
\texttt{(c) A ciascuno di voi viene chiesto di inviare una proposta segreta indicando gli oggetti che desiderate, formattata tra parentesi quadre come segue: "{[}Proposta: <numero> <nome oggetto>, <numero> <nome oggetto>, <...>{]}".} \\
\texttt{(d) Se le vostre proposte sono complementari, cioè ci sono abbastanza oggetti per soddisfare entrambe le proposte, a ciascun giocatore viene assegnato un punteggio basato sulla somma dei valori degli oggetti ricevuti. In caso contrario, entrambi ricevono zero punti.} \\
\texttt{(e) L'obiettivo è massimizzare il proprio punteggio.} \\
\\ 
\texttt{Cominciamo.} \\
\\ 
\texttt{L'insieme degli oggetti disponibili è:} \\
\\ 
\texttt{1 lampada, 3 ditali, 1 corda} \\
\\ 
\texttt{La funzione valore segreta è:} \\
\\ 
\texttt{lampada: 3, ditale: 1, corda: 4} \\
\\ 
\texttt{IMPORTANTE: i vostri messaggi, a meno che non si tratti di una proposta segreta, vengono trasmessi direttamente all'altro giocatore, quindi non includete alcuna risposta alle regole o testo di annuncio del vostro messaggio. Per fare una proposta segreta, utilizzate il formato indicato. Non utilizzare delle parentesi quadre quando si comunica all'altro giocatore, altrimenti verrà interpretata come una proposta segreta.} \\
\\ 
\\ 
\texttt{Ciao! Vedo che abbiamo a disposizione 1 lampada, 3 ditali e 1 corda da dividere tra noi. Che tipo di oggetti ti interessano di più? Io sono abbastanza flessibile e vorrei trovare una divisione che funzioni bene per entrambi.} \\
            }
        }
    }
    & \\ \\

    \theutterance \stepcounter{utterance}  
    & & & \multicolumn{4}{p{0.6\linewidth}}{
        \cellcolor[rgb]{0.9,0.9,0.9}{
            \makecell[{{p{\linewidth}}}]{
                \texttt{\tiny{[GM$\langle$P2]}}
                \texttt{Ciao! Sono interessato principalmente alla corda e alla lampada. I ditali non sono una priorità per me. Che ne dici se io prendo la corda e tu prendi la lampada? Per i ditali potremmo dividerli, magari tu ne prendi 2 e io 1, oppure se ti interessano di più puoi prenderli tutti tu. Cosa ne pensi?} \\
            }
        }
    }
     \\ \\

    \theutterance \stepcounter{utterance}  
    & & & \multicolumn{2}{p{0.3\linewidth}}{
        \cellcolor[rgb]{0.9,0.9,0.9}{
            \makecell[{{p{\linewidth}}}]{
                \texttt{\tiny{[GM$|$GM]}}
                \texttt{continue} \\
            }
        }
    }
    & & \\ \\

    \theutterance \stepcounter{utterance}  
    & & & \multicolumn{2}{p{0.3\linewidth}}{
        \cellcolor[rgb]{0.9,0.9,0.9}{
            \makecell[{{p{\linewidth}}}]{
                \texttt{\tiny{[GM$|$GM]}}
                \texttt{Ciao! Sono interessato principalmente alla corda e alla lampada. I ditali non sono una priorità per me. Che ne dici se io prendo la corda e tu prendi la lampada? Per i ditali potremmo dividerli, magari tu ne prendi 2 e io 1, oppure se ti interessano di più puoi prenderli tutti tu. Cosa ne pensi?} \\
            }
        }
    }
    & & \\ \\

    \theutterance \stepcounter{utterance}  
    & & \multicolumn{4}{p{0.6\linewidth}}{
        \cellcolor[rgb]{0.9,0.9,0.9}{
            \makecell[{{p{\linewidth}}}]{
                \texttt{\tiny{[P1$\langle$GM]}}
                \texttt{Ciao! Sono interessato principalmente alla corda e alla lampada. I ditali non sono una priorità per me. Che ne dici se io prendo la corda e tu prendi la lampada? Per i ditali potremmo dividerli, magari tu ne prendi 2 e io 1, oppure se ti interessano di più puoi prenderli tutti tu. Cosa ne pensi?} \\
            }
        }
    }
    & \\ \\

    \theutterance \stepcounter{utterance}  
    & \multicolumn{4}{p{0.6\linewidth}}{
        \cellcolor[rgb]{0.9,0.9,0.9}{
            \makecell[{{p{\linewidth}}}]{
                \texttt{\tiny{[P1$\rangle$GM]}}
                \texttt{La tua proposta mi sembra ragionevole! Dato che i ditali non sono una tua priorità e a me interessano, sarei contento di prendere la lampada e tutti e 3 i ditali, mentre tu potresti prendere la corda. In questo modo otterremmo entrambi quello che preferiamo. Ti va bene questa divisione?} \\
            }
        }
    }
    & & \\ \\

    \theutterance \stepcounter{utterance}  
    & & & \multicolumn{2}{p{0.3\linewidth}}{
        \cellcolor[rgb]{0.9,0.9,0.9}{
            \makecell[{{p{\linewidth}}}]{
                \texttt{\tiny{[GM$|$GM]}}
                \texttt{continue} \\
            }
        }
    }
    & & \\ \\

    \theutterance \stepcounter{utterance}  
    & & & \multicolumn{2}{p{0.3\linewidth}}{
        \cellcolor[rgb]{0.9,0.9,0.9}{
            \makecell[{{p{\linewidth}}}]{
                \texttt{\tiny{[GM$|$GM]}}
                \texttt{La tua proposta mi sembra ragionevole! Dato che i ditali non sono una tua priorità e a me interessano, sarei contento di prendere la lampada e tutti e 3 i ditali, mentre tu potresti prendere la corda. In questo modo otterremmo entrambi quello che preferiamo. Ti va bene questa divisione?} \\
            }
        }
    }
    & & \\ \\

    \theutterance \stepcounter{utterance}  
    & & \multicolumn{4}{p{0.6\linewidth}}{
        \cellcolor[rgb]{0.9,0.9,0.9}{
            \makecell[{{p{\linewidth}}}]{
                \texttt{\tiny{[GM$\rangle$P2]}}
                \texttt{La tua proposta mi sembra ragionevole! Dato che i ditali non sono una tua priorità e a me interessano, sarei contento di prendere la lampada e tutti e 3 i ditali, mentre tu potresti prendere la corda. In questo modo otterremmo entrambi quello che preferiamo. Ti va bene questa divisione?} \\
            }
        }
    }
    & \\ \\

    \theutterance \stepcounter{utterance}  
    & & & \multicolumn{4}{p{0.6\linewidth}}{
        \cellcolor[rgb]{0.9,0.9,0.9}{
            \makecell[{{p{\linewidth}}}]{
                \texttt{\tiny{[GM$\langle$P2]}}
                \texttt{Perfetto! Questa divisione mi sembra ideale {-} io prendo la corda che mi interessa di più, e tu prendi la lampada e tutti i ditali. Siamo d'accordo allora: tu ricevi 1 lampada e 3 ditali, io ricevo 1 corda.} \\
\\ 
\texttt{{[}Proposta: 1 corda{]}} \\
            }
        }
    }
     \\ \\

    \theutterance \stepcounter{utterance}  
    & & & \multicolumn{2}{p{0.3\linewidth}}{
        \cellcolor[rgb]{0.9,0.9,0.9}{
            \makecell[{{p{\linewidth}}}]{
                \texttt{\tiny{[GM$|$GM]}}
                \texttt{proposal} \\
            }
        }
    }
    & & \\ \\

    \theutterance \stepcounter{utterance}  
    & & & \multicolumn{2}{p{0.3\linewidth}}{
        \cellcolor[rgb]{0.9,0.9,0.9}{
            \makecell[{{p{\linewidth}}}]{
                \texttt{\tiny{[GM$|$GM]}}
                [0, 0, 1]
            }
        }
    }
    & & \\ \\

    \theutterance \stepcounter{utterance}  
    & & \multicolumn{4}{p{0.6\linewidth}}{
        \cellcolor[rgb]{0.9,0.9,0.9}{
            \makecell[{{p{\linewidth}}}]{
                \texttt{\tiny{[P1$\langle$GM]}}
                \texttt{L'altro giocatore ha presentato una proposta segreta. Ora tocca a Lei presentare una sua proposta.} \\
            }
        }
    }
    & \\ \\

    \theutterance \stepcounter{utterance}  
    & \multicolumn{4}{p{0.6\linewidth}}{
        \cellcolor[rgb]{0.9,0.9,0.9}{
            \makecell[{{p{\linewidth}}}]{
                \texttt{\tiny{[P1$\rangle$GM]}}
                \texttt{{[}Proposta: 1 lampada, 3 ditali{]}} \\
            }
        }
    }
    & & \\ \\

    \theutterance \stepcounter{utterance}  
    & & & \multicolumn{2}{p{0.3\linewidth}}{
        \cellcolor[rgb]{0.9,0.9,0.9}{
            \makecell[{{p{\linewidth}}}]{
                \texttt{\tiny{[GM$|$GM]}}
                \texttt{proposal} \\
            }
        }
    }
    & & \\ \\

    \theutterance \stepcounter{utterance}  
    & & & \multicolumn{2}{p{0.3\linewidth}}{
        \cellcolor[rgb]{0.9,0.9,0.9}{
            \makecell[{{p{\linewidth}}}]{
                \texttt{\tiny{[GM$|$GM]}}
                [1, 3, 0]
            }
        }
    }
    & & \\ \\

    \theutterance \stepcounter{utterance}  
    & & & \multicolumn{2}{p{0.3\linewidth}}{
        \cellcolor[rgb]{0.9,0.9,0.9}{
            \makecell[{{p{\linewidth}}}]{
                \texttt{\tiny{[GM$|$GM]}}
                [[1, 3, 0], [0, 0, 1]]
            }
        }
    }
    & & \\ \\

\end{supertabular}
}

\end{document}
